% !TeX root = ../thesis.tex
\chapter{Background and Related Work}
\label{chap:background}

This chapter develops the concepts needed to interpret the implemented pipeline
and its evidence. The organization is thematic: it begins with the supervision
and task distinctions created by UCF-Crime, then connects MIL, video representation learning, parameter-efficient adaptation,
prototype guidance, response-label prediction, and temporal grounding. The
human-oversight and deployment implications these ideas raise are addressed
together with the evidence in Chapter~\ref{chap:discussion}, where they can be
grounded in measured failure modes rather than anticipated in the abstract.
The aim here is not to catalogue every video anomaly method. It is to make
clear which established ideas the thesis combines, what each idea can support,
and which conclusions still require direct evaluation.

\section{Video Anomaly Tasks under Weak Supervision}
\label{sec:background-tasks}

\subsection{UCF-Crime and Long Untrimmed Video}

An anomaly is meaningful only relative to a context and an expected range of
behavior. Chandola et al. define anomaly detection as finding patterns that do
not conform to expected behavior, and they describe contextual anomalies, which
are anomalous in one context but normal in another \cite{chandola2009anomaly}.
In surveillance video, the same visual action can be routine in one
location and concerning in another. The learning problem is therefore not a
closed catalogue of every possible abnormal event. It is a decision problem
defined by a dataset, its labels, and the operating context represented by its
cameras. This dependence is especially important when model inputs include
background, viewpoint, compression, and scene layout as well as the event of
interest.

UCF-Crime was introduced to move anomaly detection toward long, untrimmed,
real-world surveillance recordings. The dataset contains 1,900 videos totaling
approximately 128 hours, with normal videos and 13 anomaly categories such as
Abuse, Arson, Burglary, RoadAccidents, Robbery, and Shooting
\cite{sultani2018real}. The authors presented two uses: binary anomaly detection,
where the anomaly categories form one positive group, and anomalous-activity
recognition, where the individual categories are distinguished. For the
anomalous test videos, the authors also provide temporal annotations, the start
and end frames of each event, obtained by averaging the labels of several
annotators. These
facets make the dataset suitable for asking different questions, but they do not
make the answers interchangeable.

The videos are untrimmed: an anomalous video can begin and end with normal
activity, and the event may occupy only a small part of its duration. A
video-level anomaly label therefore states that relevant evidence occurs
somewhere in the video; it does not state that every sampled clip is anomalous.
Similarly, attaching the video's category to a sampled clip is a weak training
approximation rather than a verified clip annotation. This distinction matters
for both representation adaptation and downstream MIL. Treating inherited clip
labels as fully observed ground truth would conceal the central uncertainty of
the task.

Video anomaly detection has been approached with different amounts of
supervision. One line of work learns from regular video alone; Hasan et al., for
example, learn a generative model of regular motion patterns with an
autoencoder \cite{hasan2016learning}. Weakly supervised methods instead train
on video-level labels of both normal and anomalous videos, as in the ranking
approach of Sultani et al. discussed in Section~\ref{sec:background-mil}.
The two settings use different benchmarks (Table~\ref{tab:datasets}). UCSD
Ped1/Ped2 \cite{mahadevan2010anomaly}, CUHK Avenue \cite{lu2013abnormal}, and
the ShanghaiTech Campus dataset \cite{liu2018future} were built for the
normal-only setting. Their anomalies are mostly local deviations such as
vehicles among pedestrians, running, or loitering. Each video comes from one of a
few fixed scenes, and the standard training split contains no
anomalies. Zhong et al.\ later reorganized ShanghaiTech into a weakly
supervised split of 238 training and 199 test videos \cite{zhong2019graph}, but
its anomaly types remain unnamed. XD-Violence is larger and weakly labeled, but
it mixes movies, sports, games, and hand-held recordings with surveillance
footage and centers on violence with audio \cite{wu2020not}.

UCF-Crime is the only one of these benchmarks that combines four properties
this thesis needs: real surveillance footage from many scenes, named incident
categories from which a response mapping can be derived, weak video-level
labels for training, and frame-level annotations of the anomalous test videos.
It is also the benchmark on which most of the weakly supervised methods compared
in Section~\ref{sec:background-comparison} report results.

\begin{table}[!ht]
    \centering
    \small
    \caption[Video anomaly datasets.]{Video anomaly datasets as described by their authors or, for UCSD, Avenue, and ShanghaiTech, by Liu et al.\ \cite{liu2018future}. ``Normal only'' means that the standard training split contains no anomalies.}
    \label{tab:datasets}
    \begin{tabularx}{\textwidth}{@{}l l >{\raggedright\arraybackslash}X l@{}}
        \toprule
        \textbf{Dataset} & \textbf{Videos (train / test)} & \textbf{Content and anomaly types} & \textbf{Training labels} \\
        \midrule
        UCSD Ped1 \cite{mahadevan2010anomaly} & 34 / 36 & Pedestrian scene; 40 events, vehicles such as bicycles and cars & Normal only \\
        UCSD Ped2 \cite{mahadevan2010anomaly} & 16 / 12 & Pedestrian scene; 12 events, same definition & Normal only \\
        CUHK Avenue \cite{lu2013abnormal} & 16 / 21 & One scene; 47 events such as throwing, loitering, running & Normal only \\
        ShanghaiTech \cite{liu2018future} & 330 / 107 & 13 scenes; 130 events of various, uncategorized types & Normal only \\
        UCF-Crime \cite{sultani2018real} & 1,610 / 290 & Real surveillance, 128\,h; 13 named crime and accident categories & Video level \\
        XD-Violence \cite{wu2020not} & 3,954 / 800 & Movies, web, and CCTV with audio, 217\,h; 6 violence types & Video level \\
        \bottomrule
    \end{tabularx}
\end{table}

A more recent direction brings in language. VadCLIP adapts a frozen CLIP
(Contrastive Language--Image Pre-training) model
with two branches, one for coarse binary classification and one that aligns
video features with textual category labels \cite{wu2024vadclip}. Anomize
targets open-vocabulary detection, in which anomaly categories not seen during
training must still be detected and named \cite{anomize2025cvpr}. UCA annotates
UCF-Crime with sentence-level event descriptions and their timing to study
surveillance video-and-language understanding \cite{yuan2024surveillance}, and, in a preprint,
Holmes-VAD fine-tunes a multimodal large language model to explain the
anomalies it detects \cite{zhang2024holmes}. In these works, understanding
means describing or naming an event in language. The pipeline in this thesis
takes a narrower route. It uses visual features only, produces no free-text
description, and predicts a fixed vocabulary of three response channels, so its
use of the term is limited to a video-level anomaly decision, response
channels, and temporal support. It also runs offline on stored videos, whereas
Karim et al. train a weakly supervised detector end to end for real-time
decisions \cite{realtime2024wacv}.

In this thesis, UCF-Crime is consequently both the empirical domain and a scope
boundary. Its video identities, fixed split, categories, and temporal annotations
define the reported questions. The dataset does not supply Police, Medical, or
Fire decisions, and it does not establish generalization to a new camera network
or institution. For the present work, this means that UCF-Crime supports weakly supervised video-level learning and test-time temporal
measurement, but response labels must be identified as a separate,
project-defined transformation.

\subsection{Detection, Classification, and Temporal Localization}

Three output levels recur throughout this thesis. \emph{Binary anomaly
detection} assigns a video, or sometimes a temporal unit within it, to Normal or
Anomaly. At video level, the output answers whether the recording contains an
event from the dataset's anomaly group. A score can be evaluated without fixing
a decision threshold through AUROC, which equals the probability that a randomly
chosen positive is ranked above a randomly chosen negative
\cite{fawcett2006introduction}. Precision, recall, and F1 instead require a
decision rule. When negatives greatly outnumber positives, a ROC curve can look
favorable while many false positives remain, and precision--recall analysis
gives the more informative picture
\cite{davis2006relationship,saito2015precision}. The distinction matters here
because the video-level test split is nearly balanced, with 140 anomalous and
150 normal videos, whereas annotated anomalous frames usually cover only part
of a video.

\emph{Anomaly classification} assigns semantic labels. The historical task in
this thesis is single-label 14-class classification: Normal plus the 13
UCF-Crime anomaly categories form mutually exclusive outcomes. The later
response task is different. It predicts three non-exclusive channels and can
activate more than one at a time. A macro-averaged F1 score (Macro-F1) from the
14-class softmax problem is therefore not commensurate with a Macro-F1 value
averaged over three
binary response channels. Label spaces, losses, thresholds, supports, and error
opportunities all differ.

\emph{Temporal localization} asks where an event occurs. It requires evidence
at a finer temporal resolution than the video label. In weakly supervised video
anomaly detection, snippet scores or pooling weights may be projected back to
time even though interval labels were not used for training. This produces a
measurable temporal signal, but one that should be evaluated explicitly against
annotations. Video-level correctness alone cannot establish accurate event
boundaries, because a model may classify a video from scene context or attend to
only a small, displaced part of the event.

These tasks form a hierarchy of questions, not a hierarchy of guaranteed
capabilities. A correct binary decision can coexist with the wrong semantic
label; a correct response vector can coexist with a temporally displaced
interval; and a localized high score can coexist with a wrong video decision.

\section{Multiple Instance Learning for Weakly Labeled Video}
\label{sec:background-mil}

\subsection{Bags, Instances, and Ambiguous Positive Evidence}

Multiple instance learning names exactly this situation: one labeled bag
contains several instances, and the instance responsible for a positive label
is unknown \cite{dietterich1997solving}. Its classical assumption is binary --
a positive bag holds at least one positive instance, a negative bag holds none
-- and modern neural MIL keeps that separation between bag-level supervision
and latent instance relevance even as it relaxes the original geometric model.

For video anomaly detection, a complete video becomes a bag and clips or
snippets become its instances. Normal videos provide negative bags, whereas an
anomalous video provides a positive bag with an unknown mixture of normal and
anomalous temporal content. Sultani et al. used this view to learn an anomaly
ranking function from video-level labels, comparing high-scoring instances from
anomalous and normal bags while adding sparsity and temporal smoothness
constraints \cite{sultani2018real}. The formulation directly addresses the
annotation problem: it does not require the anomalous segment to be labeled
during training.

Later weakly supervised approaches illustrate why instance selection and
representation quality are coupled. The Multiple Instance Self-Training
framework (MIST) alternates MIL-based pseudo-label
generation with refinement of task-specific video representations
\cite{feng2021mist}. Robust Temporal Feature Magnitude learning (RTFM) focuses on
feature magnitudes and selects discriminative
instances while modeling temporal dependencies \cite{tian2021weakly}. These
approaches differ in architecture and objective, but both address a common
difficulty: the normal snippets that dominate an anomalous video's duration can
overwhelm the comparatively rare positive evidence.

Other weakly supervised methods select or correct instance evidence in
different ways. RTFM ranks snippets by feature magnitude, and MGFN (Magnitude-Contrastive
Glance-and-Focus Network) keeps
magnitudes but adds a magnitude-contrastive loss, motivated by the observation
that feature magnitudes are inconsistent across scenes
\cite{tian2021weakly,chen2023mgfn}. Zhong et al. treat the weak labels as noisy
snippet labels, clean them with a graph convolutional network, and then train a
supervised action classifier \cite{zhong2019graph}. CLAWS (Clustering Assisted
Weakly Supervised Learning with normalcy Suppression) combines a
normalcy-suppression mechanism, which lowers the anomaly scores of normal
regions, with a clustering-distance loss that separates normal and anomalous
representations \cite{zaheer2020claws}. Li et al. use sequences of several
snippets instead of single snippets as the unit of selection and refine the
snippet scores by self-training with a transformer \cite{li2022self}. The
attention pooling described in the next subsection ranks instances by learned
weights rather than by feature magnitude, and in the binary pool every instance
keeps a non-negative weight. The VadCLIP authors also note that most current
methods operate on features precomputed by a pretrained extractor such as I3D
(Inflated 3D ConvNet)
\cite{wu2024vadclip,carreira2017quo}; the pipeline here likewise works on stored
clip features.

The bag assumption is useful but imperfect for surveillance video. A video-level
category may be uncertain, the event may be visually subtle, and a nominally
normal video cannot enumerate every acceptable behavior. MIL also does not by
itself specify how to aggregate evidence, how many instances to sample, or how
to handle a positive bag in which selected clips miss the event. Those choices
become part of the method and must be assessed under the actual sampling
protocol.

MIL therefore supplies the correct abstraction for weak video labels, but
not a guarantee of instance discovery. Fixed-size training bags, fixed-bag
evaluation, and dense variable-length evaluation expose different subsets of
the temporal evidence. Their results must therefore be named by input protocol
rather than treated as interchangeable measurements of one abstract model.

\subsection{Pooling, Attention, and Top-Instance Selection}

A MIL model requires a permutation-invariant operation that converts a set of
instance representations into a bag representation. Maximum pooling embodies a
strong positive-instance assumption: the most activated instance determines the
bag. Mean pooling distributes influence across all instances, which can suppress
a brief event when normal content dominates. Top-$k$ pooling lies between these
extremes by using a selected subset, but its behavior depends on both $k$ and the
ranking function.

Attention-based MIL learns non-negative instance weights that sum to one and
forms a weighted bag representation -- a trainable, permutation-invariant
alternative to fixed pooling operators \cite{ilse2018attention,zaheer2017deep}.
The weights identify which instances contribute most to the aggregate and
accommodate bags of different sizes, without requiring an externally specified
positive clip during training.

Attention pooling and top-instance pooling answer related but different needs.
The weighted mean gives the classifier a differentiable global summary. A
top-attention subset allows a second branch to retain high-ranked instances and
combine statistics such as their element-wise maximum and mean. Maximum values
can preserve a strong feature response, while the mean describes what is shared
across the selected subset. Concatenating these statistics with the attention
pool provides more than one view of a bag without claiming that any individual
view is sufficient.

Softmax attention is also only relative within a bag: adding instances changes
the normalizing denominator and can shift every weight even when one clip's
representation is unchanged, so a threshold applied after per-video smoothing
is not a global probability threshold. The same effect scales the top-quarter
selection rule from four clips in a 16-instance training bag to many more
clips in dense evaluation, from the identical checkpoint.

The implementation studied here uses one shared attention scorer for binary
pooling and for ranking instances used by the response branch. This reuses a
common temporal ordering, but it does not prove that the ordering is optimal
for both tasks.

\section{Video Representation Learning and Parameter-Efficient Adaptation}
\label{sec:background-representation}

\subsection{VideoMAE Representations}

Video inputs contain appearance and motion across space and time. A clip encoder
must transform that high-dimensional signal into a representation that can be
stored, sampled, and aggregated. Transformer-based video encoders divide a clip
into spatio-temporal tokens and use self-attention layers to model relations
among them. The design follows the Transformer \cite{vaswani2017attention} and
its application to images as sequences of patches in the Vision Transformer
\cite{dosovitskiy2020image}. The output can remain token-level for fine spatial
reasoning or be pooled into one vector per clip for downstream temporal
modeling.

VideoMAE learns such representations through masked autoencoding
\cite{tong2022videomae}: during self-supervised pretraining it masks a high
proportion of spatio-temporal tube tokens and reconstructs the missing content,
exploiting temporal redundancy to force the encoder toward structure that
survives reconstruction. The published model is a general video representation
learner; it does not encode the UCF-Crime response mapping or the evaluation
protocol used in this thesis.

Using a pretrained clip encoder separates two learning problems. The encoder
maps pixels to features, while the MIL model learns how to combine features
across a video. This separation can make long-video experiments tractable
because clip embeddings are extracted once and reused. It also creates a strong
provenance dependency: after features are frozen, later experiments are
conditional on the exact encoder checkpoint, preprocessing, software loader,
sampling choices, and stored feature identities.

The pipeline uses globally averaged VideoMAE token states rather than retaining
patch-level tokens. This provides a compact clip representation for MIL, but
removes the information needed for spatial localization within a frame. Dense
temporal sampling can improve coverage along the time axis; it cannot recreate
discarded spatial tokens.

\subsection{Low-Rank Adaptation}

Full fine-tuning updates every parameter of a pretrained encoder, at a
corresponding optimizer-state and storage cost. LoRA avoids that cost by
freezing the pretrained weight matrix and representing its task-specific
update as the product of two trainable low-rank matrices instead
\cite{hu2021lora}. If a frozen projection is
$W_0$, the adapted projection has the form
\begin{equation}
    W' = W_0 + \frac{\alpha}{r}BA,
\end{equation}
where $r$ is the chosen rank and $A$ and $B$ are the trainable factors. When
$r$ is much smaller than the input and output dimensions, the number of
trainable parameters is substantially smaller than for a dense update.

LoRA is one of several parameter-efficient alternatives to full fine-tuning.
Adapter modules insert small trainable layers into a frozen network
\cite{houlsby2019parameter}, and ST-Adapter, AIM, and AdaptFormer use such
modules to transfer pretrained image transformers to video or to adapt
pretrained vision transformers to new recognition tasks
\cite{pan2022st,yang2023aim,chen2022adaptformer}. Prompt tuning takes a
different route and learns a small set of input tokens while the backbone stays
frozen \cite{lester2021power,jia2022visual}. He et al. show that several of
these methods can be described as modifications of specific hidden states of
the pretrained model \cite{he2022towards}. No comparison among these
alternatives was run for this thesis; LoRA was a design choice, and its
selection is not a result.

LoRA was introduced for large language models, so applying it to a video
transformer is a transfer of an adaptation mechanism rather than a new LoRA
method. The practitioner must still choose which projections receive adapters, the rank,
scale, dropout, optimization schedule, and selection criterion. Parameter
efficiency does not imply compute-free adaptation, and a small trainable subset
does not guarantee better generalization than frozen features or full
fine-tuning.

In a sequential pipeline, adapter selection also precedes downstream evidence.
Once the adapted encoder is frozen and the feature store is extracted, the MIL
model cannot correct the encoder through pixel-level gradients. Conversely,
repeating only MIL training does not measure variation in adapter training. A
three-seed MIL result over one feature store is therefore narrower than a
three-seed repetition of the whole pipeline.

For this thesis, LoRA has a limited and defensible rationale: it provides a
parameter-efficient task-adaptation stage before frozen feature extraction. The
final results evaluate the executed LoRA--prototype--MIL sequence, but there is
no matched test experiment that isolates a causal performance gain from LoRA.
This boundary is essential to the interpretation of RQ1.

\section{Normal Prototypes and Guided Instance Sampling}
\label{sec:background-prototypes}

Distance from a representation of normality is a recurring idea in anomaly
detection. A single global center can be too restrictive when routine behavior
has several modes, while a set of prototypes can summarize multiple regions of
the normal feature space. A clip far from its nearest normal prototype may be
unusual under that representation, although distance alone does not reveal the
semantic cause or prove that the clip is anomalous.

Prototype memories have also appeared in weakly supervised video anomaly
detection. Uncertainty Regulated Dual Memory Units (UR-DMU) learns normal and
abnormal memory units and uses them to
separate feature structure under weak supervision \cite{zhou2023ur}. The present
pipeline makes a narrower construction. It clusters only embeddings from normal
training videos, then defines a clip score by squared distance to the nearest
normal center. It does not reproduce the dual-memory architecture, uncertainty
regulation, or training objective of UR-DMU.

The distance can affect learning in two conceptually separate ways. As an
appended scalar feature, it lets the MIL model decide whether distance is useful
alongside the clip embedding. As a sampling guide, it increases the chance that
high-distance clips enter a fixed-size training bag after an initial uniform
warm-up. The first mechanism changes what the model observes at each instance;
the second changes which instances the model receives. Either can help, have no
effect, or amplify irrelevant novelty such as camera changes and unusual but
normal motion.

The code calls this frozen scoring and selection role the \emph{teacher};
this thesis uses the term \emph{prototype distance}. Neither term means that the distance is calibrated ground truth,
that it was trained from temporal annotations, or that the final attention must
be monotonic in distance. The mechanisms therefore require a factorial
comparison: score input on or off, crossed with guided or uniform sampling. In
this thesis that comparison is validation-only, so it can describe behavior
within the frozen validation protocol but cannot establish a general causal or
test-set benefit.

\section{Multi-Label Response Prediction, Imbalance, and Gating}
\label{sec:background-response}

\subsection{From Exclusive Categories to Non-Exclusive Channels}

Single-label classification assumes that exactly one class is correct. A
softmax output implements this competition by normalizing scores across classes.
Multi-label classification instead permits several labels to be positive at
once. Zhang and Zhou describe it as learning from examples that are each
represented by a single instance but associated with a set of labels
\cite{zhang2014review}, and early work on multi-label scene classification
illustrates why one scene can require more than one semantic description
\cite{boutell2004learning}. Neural multi-label models commonly emit one logit per
label and apply an independent sigmoid, so activating one channel does not
mathematically suppress the others.

That formulation fits the structure of a response vector better than a
14-class softmax because Police, Medical, and Fire can co-occur. It does not
validate the labels themselves; the project-defined mapping used here is given
in Section~\ref{sec:task-formulation} and Table~\ref{tab:response-map}.

The reformulation also changes the meaning of error. A 14-class mistake chooses
the wrong exclusive category. A multi-label mistake can be an extra channel, a
missing channel, or a completely wrong vector. When common categories share a
channel, a classifier can achieve high recall by emitting that channel broadly.
An all-services predictor makes this degeneracy visible: it never misses a
positive channel but creates many unnecessary positives and cannot distinguish
among response combinations.

\subsection{Imbalance-Aware Learning and Matched Metrics}

Multi-label data often contain many more negative label positions than positive
ones. Asymmetric Loss (ASL) targets exactly this imbalance: it applies
different focusing factors to positive and negative targets and shifts
negative probabilities so that easy negatives receive little or no weight
\cite{ridnik2021asymmetric}, echoing the focal loss's down-weighting of
well-classified examples in dense object detection \cite{lin2017focal}. Such
weighting changes the training emphasis; it does not by itself choose a
calibrated decision threshold or guarantee balanced channel precision and
recall.

The implementation adds an important task-specific condition: response loss is
computed only for anomalous training bags. Normal videos train the binary branch
but do not directly train the response head toward the all-zero vector. This
separation makes it especially important to distinguish raw response output
from binary-gated output. A raw sigmoid can be positive for a normal video
because the response head was optimized for service discrimination conditional
on anomalous training examples.

Binary gating composes two predictions. If $\widehat{y}^{\mathrm{bin}}_j$ is the
binary anomaly decision and $\widehat{\vect{y}}^{\mathrm{resp}}_j$ is the
thresholded response vector, then the gated output is
\begin{equation}
    \widehat{\vect{y}}^{\mathrm{gated}}_j
    = \mathbb{1}\!\left[\widehat{y}^{\mathrm{bin}}_j=\text{Anomaly}\right]
      \widehat{\vect{y}}^{\mathrm{resp}}_j.
\end{equation}
The gate can remove response false positives on videos classified as Normal, but
it can also erase all response channels for a binary false negative. It changes
the system's error profile without changing the learned response logits.

Precision, recall, Macro-F1, Hamming loss, and exact match each expose a
different failure mode -- over-prediction, missed positives, channel-weighted
balance, aggregate error count, and complete-vector correctness, respectively
\cite{sokolova2009systematic} -- and are defined precisely in Methods
(\S\ref{sec:evaluation-paths}). A response model can have acceptable Macro-F1
but poor exact match if it frequently adds or omits one channel, so no single
summary is sufficient on its own.

Matched populations and baselines are crucial. Raw and gated metrics must use the
same 290 videos to isolate the computational effect of gating. Restricting
evaluation to ground-truth anomalous videos is an oracle-conditioned diagnostic,
not an operational pipeline. Comparing learned gated output with an
all-services rule behind the same binary gate further asks how much value comes
from channel selection rather than from anomaly detection alone.

\section{Temporal Grounding and Attention Caveats}
\label{sec:background-temporal}

\subsection{From Clip Scores to Temporal Evidence}

Temporal grounding connects a sequence of clip scores to annotated event
intervals. Its validity begins with geometry: every clip must be mapped to the
source frames it actually covers. Nominal sampling rates are not enough when
source frame rates differ or integer decimation is used. Coverage should be
reported separately from localization performance because a method cannot align
with an event it never observes, while observing every annotated frame does not
ensure that it scores those frames highly.

Several metrics expose different properties. Temporal average precision (AP)
evaluates the ordering of projected frame scores across thresholds. It is
threshold-free but does not directly measure whether a fixed predicted interval
has accurate boundaries. Top-clip hit asks whether the single highest-scoring
clip intersects any ground-truth interval; it is intuitive but gives the same
credit to a one-frame intersection and broad overlap. Interval intersection over
union (IoU) compares predicted and annotated temporal support under a declared
rule and penalizes both missed event frames and excess predicted duration.
Temporal action localization evaluates detections by their IoU with annotated
segments at several thresholds: the THUMOS challenge uses thresholds from 0.1
to 0.5 \cite{idrees2017thumos}, and W-TALC (Weakly-Supervised Temporal Activity
Localization and Classification) additionally reports 0.7
\cite{paul2018wtalc}. The thresholds 0.1, 0.3, 0.5, and 0.7 used in this thesis
are its own declared choice within this range. The frequency of exactly zero IoU keeps complete failures visible when a mean could
hide them.

These metrics can disagree. A model can rank anomalous frames above routine
frames and obtain useful AP while its smoothing and threshold rule yields broad
or displaced intervals. Conversely, predicting the whole video guarantees
non-zero overlap whenever an event exists but provides no temporal selectivity.
Frozen whole-video, uniform, random, and prototype-distance references help
interpret whether a learned attention sequence adds rank or boundary information
beyond such trivial behavior.

Scores should be aggregated per video before averaging because both event
prevalence and attention normalization vary with video length. Pooling every
frame from every video into one evaluation would allow long recordings to
dominate and would mix score scales that are only relative within their bags.
The thesis therefore uses per-video temporal metrics on a declared eligible
population and reports annotation-boundary exclusions rather than silently
clipping ground truth.

The synthesis is that temporal grounding is multi-dimensional: coverage, rank
quality, top-instance intersection, fixed-rule boundary overlap, and complete
failure frequency answer different questions. RQ4 consequently requires all of
them rather than a selected qualitative heatmap.

\subsection{Attention Is Not Calibrated Explanation}

Attention weights are useful computationally because they determine how
instances contribute to a pooled representation. It is tempting to read a large
weight as a causal explanation of the prediction. Evidence from standard
attention models shows that substantially different attention distributions can
sometimes produce similar predictions and that attention need not agree with
other feature-importance measures \cite{jain2019attention}. Although that study
examined natural-language tasks rather than surveillance video, it directly
challenges the general inference from ``high attention'' to ``faithful
explanation.''

The limitation is stronger here because the attention scorer is trained from
bag labels, not interval labels. Its primary optimization role is to form a
representation useful for video-level losses. A clip can receive weight because
it contains event evidence, because it carries correlated scene information, or
because emphasizing it helps the classifier in another way. Softmax weights also
sum to one even when no truly anomalous clip exists, so every normal video has a
highest-attended clip.

Measured overlap with temporal annotation remains informative. It tests whether
the learned ordering aligns empirically with event support more than frozen
references. The appropriate term is \emph{temporal grounding}: a relationship
between model saliency and annotated time. This does not turn the weights into
calibrated probabilities, prove that highlighted frames caused the prediction,
or establish precise event localization.

In this thesis, attention is therefore used twice but interpreted narrowly: it
is a pooling and ranking mechanism in the model, and its dense outputs are a
signal evaluated for temporal alignment. Claims of faithful explanation or
causal attribution remain outside the evidence.

\section{Comparison of Related Approaches}
\label{sec:background-comparison}

Table~\ref{tab:related-comparison} places the weakly supervised methods
discussed in this chapter side by side. The AUC column repeats the frame-level
values that each paper reports for its own method on UCF-Crime; the values were
obtained under different features and training settings and are listed to
orient the reader, not to rank the methods. Frame-level AUC has also been argued
to reward differences between recording sources as well as event localization,
and to lack the resolution to separate small margins at the size of the
UCF-Crime test split \cite{abdulaziz2026auditing}. The video-level AUROC reported in this thesis is
not comparable with these values (Section~\ref{sec:results-protocol}).

\begin{table}[!ht]
    \centering
    \small
    \caption[Related weakly supervised approaches on UCF-Crime.]{Related weakly supervised approaches on UCF-Crime. AUC is the frame-level value each paper reports for its own method; the values use different features and protocols and are not a ranking. Every method outputs a binary anomaly score; ``adds'' marks further outputs.}
    \label{tab:related-comparison}
    \begin{tabularx}{\textwidth}{@{}l l >{\raggedright\arraybackslash}X r@{}}
        \toprule
        \textbf{Method} & \textbf{Features} & \textbf{Key idea} & \textbf{AUC (\%)} \\
        \midrule
        Sultani et al.\ \cite{sultani2018real} & C3D & MIL ranking of the top segment & 75.41 \\
        Zhong et al.\ \cite{zhong2019graph} & TSN & Graph-based label-noise cleaning & 82.12 \\
        CLAWS \cite{zaheer2020claws} & C3D & Normalcy suppression, clustering loss & 83.03 \\
        MIST \cite{feng2021mist} & I3D & Pseudo-labels, self-guided attention & 82.30 \\
        RTFM \cite{tian2021weakly} & I3D & Top-$k$ feature magnitude & 84.30 \\
        MSL \cite{li2022self} & VideoSwin & Multi-sequence top-$k$, self-training & 85.62 \\
        UR-DMU \cite{zhou2023ur} & I3D & Dual memory units, uncertainty & 86.97 \\
        MGFN \cite{chen2023mgfn} & I3D & Glance-and-focus, magnitude contrast & 86.98 \\
        PEL4VAD$^{*}$ \cite{pu2023learning} & I3D & Context aggregation, text prompts & 86.76 \\
        Karim et al.\ \cite{realtime2024wacv} & End-to-end & Trained jointly for real-time decisions & 86.94$^{\dagger}$ \\
        VadCLIP \cite{wu2024vadclip} & CLIP & Dual branch, text alignment; adds category & 88.02 \\
        Holmes-VAD$^{*}$ \cite{zhang2024holmes} & Multimodal LLM & Instruction-tuned explanation; adds text & 89.51 \\
        \midrule
        This thesis & VideoMAE + LoRA & Shared attention, top-quarter pooling; adds three response channels and intervals & n/a \\
        \bottomrule
    \end{tabularx}

    \smallskip
    {\footnotesize $^{*}$ Preprint. \quad $^{\dagger}$ Value taken from the paper's abstract; all other values come from the results tables.\par}
\end{table}

Three differences follow from the table. First, the compared methods output
an anomaly score, in two cases with an anomaly category or a text explanation,
while this thesis predicts three non-exclusive response channels that none of
them addresses. Second, most compared papers evaluate temporal evidence through
the pooled frame-level AUC; VadCLIP additionally reports category-level
detection precision at interval-overlap thresholds. This thesis instead
measures per-video interval overlap and average precision and compares both
with trivial temporal references. Third, the table lists one number per method,
whereas this thesis repeats Stage-2 training over three seeds and compares every
learned output with a matched trivial rule. The thesis does not propose a new detector and does not compete
on frame-level AUC. Its contribution lies in this combination of outputs and
evaluation, which Chapter~\ref{chap:methods} specifies.
